\documentclass{article}
\usepackage{amsmath}
\usepackage{amssymb}

\title{MAT 5210 Homework 7}
\author{}
\date{}

\begin{document}

\maketitle

\begin{enumerate}
\item Let $L:K$ be finite Galois and let $P$ be unramified. Show that $P$ splits completely
in $L$ if and only if $\operatorname{Frob}_{\mathfrak P}=id$ for one, and hence every,
$\mathfrak P$ above $P$.

\item Let $L=\mathbb Q(\zeta_7)$. For each nonzero congruence class of $p$ modulo $7$,
determine the residue degree $f$ and the number $g$ of prime ideals lying above $p$.
Hint: $f=\operatorname{ord}_7(p)$ and $fg=6$.

\item Let $K=\mathbb Q(\sqrt5)$. For an odd prime $p\neq5$, show that
\[
p\text{ splits in }K
\quad\Longleftrightarrow\quad
\left(\frac5p\right)=1,
\]
and interpret the two possible values of the Legendre symbol in terms of Frobenius.

\item Suppose $C_K\cong\mathbb Z_6$. Determine all possible degrees $[M:K]$ for
$K\subseteq M\subseteq H_K$, and describe the corresponding subgroups of $C_K$.

\item We computed earlier that $C_{\mathbb Q(\sqrt{-5})}\cong\mathbb Z_2$.
Without explicitly constructing $H_K$, determine $[H_K:K]$, its Galois group, and the
criterion for a prime ideal $P$ to split completely in $H_K$.


\item Let $L=\mathbb Q(i)$. For every odd positive integer $n$, show directly that
\[
\Theta_{L/\mathbb Q}((n))=
\begin{cases}
1,&n\equiv1\pmod4,\\
c,&n\equiv3\pmod4,
\end{cases}
\]
where $c(i)=-i$. Explain why this shows that principal ideals need not lie in the kernel of
the Artin map for a general abelian extension.

\item Prove directly that
\[
C_{\mathbb Q}(m\infty)\cong(\mathbb Z/m\mathbb Z)^*.
\]
Explain why the positivity condition at the real place is needed to distinguish $a$ from
$-a$.
\end{enumerate}

\end{document}
